% titel: Blatt 5, A3: Valenzelektronen in Iodid (Diskussion)
% skript: Mortimer Kap. 7 (chemie: Valenzelektronen, Edelgaskonfiguration)

Valenzelektronen sind die an chemischen Bindungen beteiligten Elektronen der
äussersten, unvollständig besetzten Schale. Für Hauptgruppenelemente sind das
die $ns$- und $np$-Elektronen, für Übergangsmetalle oft zusätzlich $(n-1)d$.

\begin{werkzeugkasten}[Wann zählt ein Elektron zu den Valenzelektronen?]
\begin{itemize}
  \item Abgrenzung über die \textit{Hauptquantenzahl $n$}: nur Elektronen der
        höchsten, im Element besetzten Schale zählen als Valenzelektronen.
  \item Abgeschlossene, energetisch tief liegende Unterschalen (z.\ B.
        $3d^{10}$, $4d^{10}$) gehören zur Rumpfkonfiguration und nehmen an der
        Bindung nicht teil.
  \item Im Anion wird die Valenzschale zur Edelgaskonfiguration ergänzt; die
        zugesetzten Elektronen landen in der höchsten Schale.
\end{itemize}
So ist z.\ B. für ein Iod-Atom die Valenzschale $5s^2\,5p^5$; im $I^-$ wird sie
zu $5s^2\,5p^6$ aufgefüllt. Die voll besetzte $4d^{10}$-Unterschale bleibt
Rumpf.
\textit{Orientierung:} PSE, Spalte VII; Iod liegt in Periode 5.
\end{werkzeugkasten}

\begin{achtung}[Darauf achten]
\begin{itemize}
  \item Zählen Sie beim Anion die Elektronen der \textit{letzten belegten
        Schale}, nicht alle Schalen unterhalb.
  \item Eine volle Unterschale $(n-1)d^{10}$ ist nicht automatisch Valenz ---
        entscheidend ist die Hauptquantenzahl der äussersten Schale.
  \item Bei $\mathrm{Br}$ und $\mathrm{I}$ (Periode 4/5) gibt es zusätzlich
        besetzbare $d$-Orbitale der gleichen Periode ($5d$), die erst bei
        Oktetterweiterung (z.\ B. in $\mathrm{IF_7}$) bindend werden --- aber die
        $4d$ von Iod bleiben voll und inert.
\end{itemize}
\end{achtung}
